\documentclass[12pt]{article}
\usepackage[utf8]{inputenc}
\usepackage[T1]{fontenc}
\usepackage{hyperref}
\title{Khlula Language Resources for K-12 Curriculum Development}
\author{Compiled by Claude Sonnet 5.5 (claude-sonnet-5-5) for Andr\'as Kornai}
\date{}
\begin{document}
\maketitle

\section{Introduction}
Khlula (ISO 639-3 ykl; also Tula, Namupha, Alapha, Black Phula) is a Southeastern Loloish language of the Phula group, spoken in Wenshan Prefecture (Wenshan and Maguan counties), Yunnan, China. Wikipedia gives about 21,000 speakers (2011) and Joshua Project 35,000. Wikitongues rates it 6b (Threatened). Mandarin is increasingly used, especially by younger people and in schools. No educational materials turned up in the search.

\section{Primary Educational Resources}
None found.

\section{Linguistic and Reference Materials}
\begin{itemize}
\item Khlula language (Wikipedia): \url{https://en.wikipedia.org/wiki/Khlula_language}
\item Zokhuo language (Wikipedia), its closest relative: \url{https://en.wikipedia.org/wiki/Zokhuo_language}
\item Endangered Languages Project bibliography: \url{https://endangeredlanguages.com/lang/4964/bibliography?hl=es}
\item Glottolog map page: \url{https://glottolog.org/resource/languoid/id/khlu1236.bigmap.html}
\end{itemize}

\section{Religious and Community Resources}
\begin{itemize}
\item Joshua Project, Khlula in China: \url{https://joshuaproject.net/people_groups/21075}
\end{itemize}

\section{Digital and Audio Resources}
\begin{itemize}
\item Wikitongues page for Khlula (contents were not checked): \url{https://wikitongues.org/languages/ykl/}
\end{itemize}

\section{Curriculum Development Context}
The search found no orthography, primers or school programmes. The language is a Mandarin-dominated minority language, so curriculum work would need new documentation.
\end{document}